
\subsection{Conservative TCP Loopback Evaluation}

We evaluate the semi-honest RM-VOLE prototype on a single host using TCP loopback.  Each data point reports the median of three repetitions.  The RM-VOLE measurement includes three components: a real libOTe noisy-VOLE split-input layer for the $2t$ sparse coefficients, a vector-valued \texttt{RegularPprf} setup that distributes the resulting payload shares, and local WHT-domain expansion by the two parties.  The measured relation is
\[
  Z_0 + Z_1 = \Delta \cdot x,
\]
where $x \in R_p$ and the output coordinates lie in a fixed-prime vector representation of $R_{p^m}$.

The external baseline is libOTe's generic noisy subfield VOLE harness, measured over the same TCP loopback transport and normalized with \(\mathrm{entries}=mN\).  This baseline realizes the analogous relation $a=b+c\Delta$, but the field representation is not identical to the RM-VOLE fixed-prime coordinates.  We therefore treat it as a generic external noisy subfield-VOLE baseline, not as a same-field or state-of-the-art comparison.  The silent subfield VOLE audit also indicates that a more appropriate silent baseline remains future work in this branch.

Table~\ref{tab:rmvole-main-conservative} summarizes the conservative comparison.  With the real noisy-VOLE split-input setup included, RM-VOLE reduces communication by 3.8--7.9$\times$ over this generic noisy baseline.  Runtime is mixed for the smallest setting: the noisy subfield baseline is faster at $N=2^{12},m=8$ and near parity at $N=2^{12},m=16$.  On the medium and larger tested settings, RM-VOLE is faster by 1.3--2.9$\times$.

\begin{table}[t]
\centering
\caption{Conservative TCP-loopback comparison.  RM-VOLE includes noisy-VOLE split-input product sharing, vector-valued PPRF setup, and local expansion.}
\label{tab:rmvole-main-conservative}
\begin{tabular}{rrrrrrrr}
\toprule
$\log_2 N$ & $t$ & $m$ & RM s & noisy subfield s & RM MB & noisy subfield MB & comm ratio \\
\midrule
12 & 8 & 8 & 0.045 & 0.020 & 1.66 & 6.30 & 3.8$\times$ \\
12 & 8 & 16 & 0.092 & 0.086 & 6.45 & 25.18 & 3.9$\times$ \\
14 & 16 & 8 & 0.067 & 0.090 & 3.28 & 25.17 & 7.7$\times$ \\
14 & 16 & 16 & 0.157 & 0.212 & 12.80 & 100.68 & 7.9$\times$ \\
14 & 16 & 32 & 0.255 & 0.513 & 50.70 & 402.68 & 7.9$\times$ \\
16 & 64 & 8 & 0.212 & 0.269 & 12.98 & 100.67 & 7.8$\times$ \\
16 & 64 & 16 & 0.306 & 0.903 & 50.87 & 402.67 & 7.9$\times$ \\
\bottomrule
\end{tabular}
\end{table}

The breakdown in Table~\ref{tab:rmvole-breakdown-conservative} explains where the RM-VOLE cost currently comes from.  The split-input noisy VOLE accounts for almost all communication, roughly 97--99\% across the measured settings, and a large fraction of runtime.  This is expected: the current split-input layer is deliberately conservative and uses noisy VOLE for only the sparse coefficient products.  It is not the intended final silent or batched sparse-product layer.  Consequently, optimizing this layer is the next target before making stronger performance claims.

\begin{table}[t]
\centering
\caption{RM-VOLE component breakdown for the conservative setup.}
\label{tab:rmvole-breakdown-conservative}
\begin{tabular}{rrrrrrrr}
\toprule
$\log_2 N$ & $t$ & $m$ & total s & split s & PPRF s & expand s & split comm \\
\midrule
12 & 8 & 8 & 0.045 & 0.034 & 0.011 & 0.001 & 97.6\% \\
12 & 8 & 16 & 0.092 & 0.079 & 0.016 & 0.001 & 99.2\% \\
14 & 16 & 8 & 0.067 & 0.039 & 0.026 & 0.004 & 97.4\% \\
14 & 16 & 16 & 0.157 & 0.111 & 0.047 & 0.009 & 99.1\% \\
14 & 16 & 32 & 0.255 & 0.215 & 0.053 & 0.015 & 99.7\% \\
16 & 64 & 8 & 0.212 & 0.092 & 0.119 & 0.017 & 97.4\% \\
16 & 64 & 16 & 0.306 & 0.185 & 0.132 & 0.029 & 99.1\% \\
\bottomrule
\end{tabular}
\end{table}

For downstream matrix-style applications, we use a simple linear rank-one cost model.  If a workload needs $q$ independent RM-VOLE rank-one instances at fixed $(N,m)$, we estimate
\[
  T_{\mathrm{matrix}}(q,N,m) = q \cdot T_{\mathrm{rank1}}(N,m),\qquad
  C_{\mathrm{matrix}}(q,N,m) = q \cdot C_{\mathrm{rank1}}(N,m).
\]
This model intentionally does not assume cross-instance batching or amortization.  It is therefore conservative for future implementations that batch split-input setup, reuse base OT material, or replace noisy VOLE with a silent sparse-product layer.

All numbers in this section are semi-honest only.  They should not be read as malicious-secure timings, and they should not be used as a headline state-of-the-art claim.  The main conclusion is narrower: after including a real noisy-VOLE split-input setup, this RM-VOLE prototype still substantially reduces communication relative to the generic noisy subfield-VOLE baseline, while runtime improves on the medium and larger tested TCP-loopback settings.
